\documentclass[a4paper]{article}

%% Language and font encodings
\usepackage[spanish]{babel}
\usepackage[utf8x]{inputenc}
\usepackage[T1]{fontenc}

%% Sets page size and margins
\usepackage[a4paper,top=3cm,bottom=2cm,left=2cm,right=2cm,marginparwidth=1.75cm]{geometry}

%% Useful packages
\usepackage{amsmath}
\usepackage{graphicx}
\usepackage{caption}
\usepackage{subcaption}
\usepackage{float}
\usepackage{titling} 
\usepackage{booktabs} % Tablas

%% Referencias/Bibliografía
%\usepackage[backend=biber, style=apa]{biblatex}
%\addbibresource{ref.bib} % Nombre del archivo con las referencias

%Links de referencias
\usepackage[hidelinks]{hyperref} % Hace que las referencias sean clicables
\newcommand{\myref}[1]{\textsuperscript{\textbf{\ref{#1}}}} % Define un estilo propio

% Código
\usepackage{xcolor}
\usepackage{listings}

% Definición de colores estilo RStudio
\definecolor{commentcolor}{RGB}{34,139,34}   % Verde
\definecolor{keywordcolor}{RGB}{0,0,255}    % Azul
\definecolor{stringcolor}{RGB}{163,21,21}     % Rojo oscuro
\definecolor{backcolor}{RGB}{245,245,245}    % Gris muy claro

\lstset{
    language=R,                     % Lenguaje de programación
    basicstyle=\ttfamily\small,      % Fuente monoespaciada pequeña
    backgroundcolor=\color{backcolor}, % Color de fondo
    commentstyle=\color{commentcolor}, % Color de comentarios
    keywordstyle=\color{keywordcolor}\bfseries, % Color de palabras clave
    stringstyle=\color{stringcolor},   % Color de cadenas de texto
    numbers=left,                    % Números de línea a la izquierda
    numberstyle=\tiny\color{gray},   % Estilo de los números
    stepnumber=1,                    % Número en cada línea
    breaklines=true,                 % Ajuste automático de línea
    frame=single,                    % Recuadro simple alrededor del código
    rulecolor=\color{black!10},      % Color del borde del recuadro
    showstringspaces=false,          % No mostrar espacios en las cadenas
    captionpos=b,                     % Título debajo del código
    literate= {á}{{\'a}}1 {é}{{\'e}}1 {í}{{\'i}}1 {ó}{{\'o}}1 {ú}{{\'u}}1 {ñ}{{\~n}}1 {Á}{{\'A}}1 {É}{{\'E}}1 {Í}{{\'I}}1 {Ó}{{\'O}}1 {Ú}{{\'U}}1 {Ñ}{{\~N}}1 % Para las tildes y caracteres del español
}

%%%%%%%% DOCUMENT %%%%%%%%
\begin{document}

%%%% Title Page
\begin{titlepage}

\newcommand{\HRule}{\rule{\linewidth}{0.5mm}} 							% horizontal line and its thickness
\center 
 
% University
\textsc{\LARGE Universidad de Granada}\\[1cm]

% Document info
\textsc{\Large Periféricos y Dispositivos de Interfaz Humana}\\[0.2cm]
\textsc{\large Curso 2025/2026}\\[1cm] 										% Course Code
\HRule \\[0.8cm]
{ \huge \bfseries Práctica 5. Experimentación con el sistema de 
salida de sonido}\\[0.7cm]								% Assignment
\HRule \\[2cm]
\large
\emph{Autora:}\\
Inés Ruiz Sánchez\\[1.5cm]													% Author info
{\large \today}\\[5cm]

\includegraphics[width=0.3\textwidth]{images/UGRlogo.png}\\[1cm] 	% University logo
\vfill 
\end{titlepage}

%%\begin{abstract}
%%Your abstract.
%%\end{abstract}

% Generar índice de contenidos y figuras
\tableofcontents
%\newpage
\listoffigures
\newpage

%%%% SECTIONS
\section{Introducción y objetivos}
Los objetivos concretos de esta práctica son: 
\begin{itemize}
    \item Identificar y representar gráficamente la forma de onda de señales de sonido.
    \item Conocer la estructura de un fichero típico de sonido (ficheros WAV). 
    \item Entender y operar con los parámetros principales de una señal de sonido.
\end{itemize}


En el anexo\myref{anexo:an} se incluyen los ficheros de códigos al completo. \\

Para obtener los ficheros WAV he usado la página \url{https://speechgen.io/}.\\

Repositorio en GitHub: \url{https://github.com/ruiz314/PDIH/tree/main/P5}.\\

Para poder trabajar con R y RStudio en mi ordenador personal realicé los seminarios y tutoriales proporcionados por el profesor. Se encuentran en el siguiente repositorio: \url{https://github.com/ruiz314/PDIH/tree/main/S-sonido}.

\section{Requisitos mínimos}
 Para superar esta práctica se proponen los siguientes ejercicios, que se podrán realizar en un mismo programa:

\begin{itemize}
    \item Crear dos ficheros de sonido (WAV) para realizar los siguientes ejercicios. En el primero debe escucharse el nombre de la persona que realiza la práctica. En el segundo debe escucharse el apellido.
    \item Leer los dos ficheros de sonido creados y dibujar la forma de onda de ambos sonidos (por separado).

    \item Obtener la información de las cabeceras de ambos sonidos.

    \item Unir ambos sonidos en uno nuevo para escuchar el nombre y apellido correctamente.

    \item Dibujar la forma de onda de la señal y reproducir el sonido resultante (una vez unidos).

    \item Almacenar el sonido resultante en un archivo nuevo llamado ``basico.wav''  
\end{itemize}

\subsection{Ejercicio 1}\label{subsec:ej1}
% Ejercicio 1.
Para crear los ficheros de sonido en la página \url{https://speechgen.io/} hay que seleccionar el idioma y la voz (arriba a la izquierda). A continuación hay que escribir el texto en el recuadro. Observar Figura \ref{fig:img1}.\\ 

\begin{figure}[htbp]
    \centering
    \includegraphics[width=0.9\linewidth]{images/1_1convertir_nombre1.png}
    \caption{Convertir texto a sonido}
    \label{fig:img1}
\end{figure}

Abajo a la izquierda hay un botón de ``SETTINGS'' (Figura \ref{fig:img2}) que sirve para cambiar la configuración del archivo.

\begin{figure}[H]
    \centering
    \includegraphics[width=0.9\linewidth]{images/1_2convertir_ajustes.png}
    \caption{Convertir texto a sonido}
    \label{fig:img2}
\end{figure}

Seleccionar formato ``wav''.\\

Por último clikar en ``CONVERT TO SPEECH''. Aparecerá abajo el audio, y podemos descargarlo pulsando el botón ``DOWNLOAD''. Observar Figura \ref{fig:img3}.

\begin{figure}[H]
    \centering
    \includegraphics[width=0.9\linewidth]{images/1_3descargar.png}
    \caption{Descargar sonido}
    \label{fig:img3}
\end{figure}

Para crear el fichero de sonido con los apellidos hay que realizar el mismo procedimiento. Observar Figura \ref{fig:img4}.
\begin{figure}[H]
    \centering
    \includegraphics[width=0.9\linewidth]{images/1_4apellidos.png}
    \caption{Fichero de sonido con apellido}
    \label{fig:img4}
\end{figure}


%----------------------------------------------------------------

\subsection{Ejercicio 2}\label{subsec:ej2}
% Ejercicio 2
Para cargar los ficheros de audio se usa la función \texttt{readWave()}:
\begin{lstlisting}[caption={Cargar los ficheros de sonido en R}, label={lst:cargar_ficheros}]
# Audio con el nombre
nombre <- readWave('nombre.wav')
nombre

# Audio con los apellidos
apellidos  <- readWave('apellidos.wav')
apellidos
\end{lstlisting}

Llamar al fichero de sonido (líneas 3 y 7) hace que se muestre información como el número de muestras, la duración y el número de canales. Observar Figura \ref{fig:img2_1}.\\

\begin{figure}[H]
    \centering
    \includegraphics[width=0.95\linewidth]{images/2_cargar.png}
    \caption{Cargar y leer los ficheros de sonido}
    \label{fig:img2_1}
\end{figure}

Los datos obtenidos son los siguientes (ver Cuadro \ref{tab:tabla1}):
\begin{table}[ht]
    \centering
    \caption{Información de los objetos Wave: \texttt{nombre} y \texttt{apellidos}}
    \label{tab:tabla1}
    \begin{tabular}{lll}
        \toprule
        \textbf{Parámetro} & \textbf{Objeto: \texttt{nombre}} & \textbf{Objeto: \texttt{apellidos}} \\
        \midrule
        Número de muestras & 32879 & 49007 \\
        Duración (segundos) & 0.68 & 1.02 \\
        Tasa de muestreo (Hz) & 48000 & 48000 \\
        Canales & Stereo & Stereo \\
        Formato PCM & TRUE & TRUE \\
        Bit & 16 & 16 \\
        \bottomrule
    \end{tabular}
\end{table}

Para obtener el dibujo de la onda la función \texttt{plot} necesita saber el número de muestras que tiene el fichero. Observamos que el fichero \textit{nombre} tiene $32879$ muestras y el fichero \textit{apellidos} tiene $49007$. 

\begin{lstlisting}[caption={Comandos para dibujar las ondas}, label={lst:fun_plot1}]
# Mostrar la onda del sonido
plot( extractWave(nombre, from = 1, to = 32879) )

plot( extractWave(apellidos, from = 1, to = 49007) )
\end{lstlisting}

Al usar la función \texttt{plot} obtenemos la gráfica que se muestra en la Figura \ref{fig:img2_2}:
\begin{figure}[H]
    \centering
    \includegraphics[width=0.95\linewidth]{images/2_2onda.png}
    \caption{Cargar y leer los ficheros de sonido}
    \label{fig:img2_2}
\end{figure}

El dibujo de la onda del archivo de sonido ``nombre.wav'' es el que se muestra en la Figura \ref{fig:img2_3}. El dibujo obtenido para el sonido de ``apellidos.wav'' se observa en la Figura \ref{fig:img2_4}.

\begin{figure}[H]
    \centering
    \includegraphics[width=0.95\linewidth]{images/2_3plot_nombre.png}
    \caption{Dibujo de fichero ``nombre.wav''}
    \label{fig:img2_3}
\end{figure}

\begin{figure}[H]
    \centering
    \includegraphics[width=0.95\linewidth]{images/2_4plot_apellidos.png}
    \caption{Dibujo de fichero ``apellidos.wav''}
    \label{fig:img2_4}
\end{figure}


%----------------------------------------------------------------

\subsection{Ejercicio 3}\label{subsec:ej3}
% Ejercicio 3. Obtener la información de las cabeceras de ambos sonidos. 
Usando la función \texttt{str()} se puede inspeccionar las cabeceras de archivos de sonido.
\begin{lstlisting}[caption={Funciones para las cabeceras de los ficheros de sonido}, label={lst:cabeceras_ficheros}]
str(nombre)
str(apellidos)    
\end{lstlisting}

En la siguiente imagen se ve la información obtenida (Figura \ref{fig:cabeceras_fich}).
\begin{figure}[H]
    \centering
    \includegraphics[width=0.95\linewidth]{images/3_cabeceras.png}
    \caption{Cabeceras de los ficheros}
    \label{fig:cabeceras_fich}
\end{figure}

Como ambos sonidos tienen dos canales, el atributo \textit{left} y el atributo \textit{right} valen lo mismo para cada sonido. Además, confirmamos que tienen dos canales porque la variable \textit{stereo} es verdadera.


%----------------------------------------------------------------

\subsection{Ejercicio 4}\label{subsec:ej4}
% Ejercicio 4. Unir ambos sonidos en uno nuevo para escuchar el nombre y apellido correctamente.
Con la función \texttt{pastew()} se unen dos sonidos. Hay que prestar atención al orden, pues el primer argumento será el segundo sonido que se escuche y el segundo argumento será el primer sonido que se escuche.
\begin{lstlisting}[caption={Función para unir dos sonidos}, label={lst:union_sonidos}]
union <- pastew(apellidos, nombre, output="Wave")
\end{lstlisting}

En la Figura \ref{fig:img4_1} se puede observar la información del sonido.
\begin{figure}[H]
    \centering
    \includegraphics[width=0.95\linewidth]{images/4_1union.png}
    \caption{Unión de dos sonidos}
    \label{fig:img4_1}
\end{figure}

En el Cuadro \ref{tab:tabla2} se pueden observar los datos de los tres ficheros. La duración del fichero \textit{union} es la suma de las duraciones del fichero \textit{nombre} y el fichero \textit{apellidos}, aproximadamente. Sucede lo mismo con el número de muestras: $32,879 + 49,007 = 81,886$.
\begin{table}[ht]
    \centering
    \caption{Comparativa de los objetos Wave: \texttt{nombre}, \texttt{apellidos} y \texttt{union}}
    \label{tab:tabla2}
    \begin{tabular}{llll}
        \toprule
        \textbf{Parámetro} & \textbf{Objeto: \texttt{nombre}} & \textbf{Objeto: \texttt{apellidos}} & \textbf{Objeto: \texttt{union}}\\
        \midrule
        Número de muestras & 32879 & 49007 & 81886\\
        Duración (segundos) & 0.68 & 1.02 & 1.71 \\
        Tasa de muestreo (Hz) & 48000 & 48000 & 48000\\
        Canales & Stereo & Stereo & Mono\\
        Formato PCM & TRUE & TRUE & TRUE\\
        Bit & 16 & 16 & 16\\
        \bottomrule
    \end{tabular}
\end{table}

%----------------------------------------------------------------

\subsection{Ejercicio 5}\label{subsec:ej5}
% Ejercicio 5. Dibujar la forma de onda de la señal y reproducir el sonido resultante (una vez unidos). 
Llamando a la función \textit{plot()} se genera una imagen con la onda resultante (Observar la Figura \ref{fig:img5}).
\begin{lstlisting}[caption={Generación de gráfica de unión}, label={lst:graf_union}]
# Generar y mostrar dibujo de la onda
plot(extractWave(union, from=1, to=81886))

# Reproducir sonido
listen(union)
\end{lstlisting}

Además, para reproducir el sonido nuevo generado se usa la función \textit{listen(union)} y abre el reproductor en el ordenador. Se puede observar cómo se abre en la Figura \ref{fig:img4_1}.
\begin{figure}[H]
    \centering
    \includegraphics[width=0.95\linewidth]{images/5_ondaUnion.png}
    \caption{Gráfica de la onda unida}
    \label{fig:img5}
\end{figure}


%----------------------------------------------------------------

\subsection{Ejercicio 6}\label{subsec:ej6}
% Ejercicio 6
Para guardar en disco el nuevo sonido generado se llama a la función \textit{writeWave()} indicándole el sonido, y el nombre del fichero:
\begin{lstlisting}[caption={Almacenar sonido union en fichero ``basico.wav''}, label={lst:almacen_union}]
writeWave(union, file.path("basico.wav"))
basico <- readWave('basico.wav')
basico
\end{lstlisting}

A continuación, en la Figura \ref{fig:img6} se puede observar que la llamada a la función no genera ningún error y que el sonido ha sido almacenado con éxito. Además, se muestran los datos de la onda y coinciden con los datos de la onda \textit{union}. 
\begin{figure}[H]
    \centering
    \includegraphics[width=0.95\linewidth]{images/6_basico.png}
    \caption{Creación de ``basico.wav''}
    \label{fig:img6}
\end{figure}


%----------------------------------------------------------------
\newpage
%% Section 2
\section{Requisitos ampliados}
Como requisitos ampliados (opcionales para subir nota) se propone:
\begin{itemize}
    \item Pasarle un filtro de frecuencia para eliminar las frecuencias entre 10.000Hz y 20.000Hz. Almacenar la señal obtenida como un fichero WAV denominado ``filtrado.wav''  
    \item Tomar el sonido que se creó antes (lo tendremos en el archivo llamado ``basico.wav'') para aplicarle el efecto de eco. Guardar ese sonido en un archivo nuevo llamado ``eco.wav''. A continuación, se le debe dar la vuelta al sonido y almacenarlo como un fichero llamado ``alreves.wav''
\end{itemize}

% Requisitos ampliados

\subsection{Ejercicio 7}\label{subsec:ej7}
La función \texttt{bwfilter()} aplica un filtro de frecuencia, eliminando las 
muestras entre las frecuencias (en Hz) indicadas.
\begin{lstlisting}[caption={Filtro de frecuencia}, label={lst:filtro}]
filtro <- bwfilter(basico, f=basico@samp.rate, from=10000, to=20000, bandpass= FALSE, output = "Wave")

writeWave(filtro, file.path("filtrado.wav"))
\end{lstlisting}

Al poner \texttt{f = basico@samp.rate} la función lee automáticamente la cabecera del archivo y aplica el filtro de frecuencias de manera exacta y sin errores.

\begin{figure}[ht]
    \centering
    \includegraphics[width=0.95\linewidth]{images/7_filtro.png}
    \caption{Filtro de frecuencia}
    \label{fig:filtro}
\end{figure}

En la Figura \ref{fig:filtro} da un aviso al ejecutar la función \texttt{writeWave}. Este aviso no afecta al ejercicio ya que el redondeo no es perceptible para el oído humano.\\

Para ver información sobre el sonido uso la función \texttt{str} y la propia llamada al sonido. Se puede observar en \textit{Listing} \ref{lst:filtroinfo} y \textit{Figura} \ref{fig:filtroinfo}.
\begin{lstlisting}[caption={Información del filtro de frecuencia}, label={lst:filtroinfo}]
str(filtro)
filtro
\end{lstlisting}

\begin{figure}[ht]
    \centering
    \includegraphics[width=0.95\linewidth]{images/7_2infoFiltro.png}
    \caption{Información del \textit{filtro}}
    \label{fig:filtroinfo}
\end{figure}
%----------------------------------------------------------------

\subsection{Ejercicio 8}\label{subsec:ej8}


%----------------------------------------------------------------

\newpage
\section{Anexo}\label{anexo:an}
% --- Requisitos mínimos (Ejer 1-6)
\subsection{Fichero con requisitos mínimos} \label{anexo:ej_obli} 
\lstinputlisting[language=R]{src/ejercicios.R}

% --- Requisitos ampliados (Ejer 7-8)
\subsection{Fichero con requisitos ampliados} \label{anexo:ej_optat} 
\lstinputlisting[language=R]{src/extra.R}


%\section{Bibliografía y referencias}
%\printbibliography

\end{document}
